\documentclass[11pt]{article}
\usepackage[utf8]{inputenc}
\usepackage{amsmath}
\usepackage{amsfonts}
\usepackage{amssymb}
\usepackage{graphicx}
\usepackage[super]{nth}
\usepackage{amsthm}
\usepackage{bm}
\newtheorem{theorem}{Theorem}
\newtheorem{objective}{Objective}
\newtheorem{model}{Model}
\usepackage{xcolor}
\definecolor{light-gray}{gray}{0.95}
\newcommand{\code}[1]{\colorbox{light-gray}{\texttt{#1}}}
\usepackage{listings}
\usepackage{placeins}
\usepackage{enumitem}
\usepackage[T1]{fontenc}

\usepackage[authoryear]{natbib}

\lstset{language=R}


\makeatletter
\renewcommand{\maketitle}{
\begin{center}

\pagestyle{empty}

{\Large \bf \@title\par}
\vspace{1cm}

{\LARGE Marcel Gietzmann-Sanders}\\[1cm]

STAT651 - Statistical Theory I \\
University of Alaska Fairbanks


\end{center}
}\makeatother


\title{Homework 9}

\date{2025}
\setcounter{tocdepth}{2}
\begin{document}
\maketitle


\section*{Question 1}

\subsection*{a)}

The function $g(x) = \sqrt{x}$ is concave and therefore:

$$g(\mathbb{E}(x)) \geq \mathbb{E}(g(x))$$

Therefore $g(\mathbb{E}(x))$ is an upper bound for $\mathbb{E}(g(x))$. As $\mathbb{E}(x) = 2.3(4.1) = 9.43$ our upper bound is $\sqrt{9.43} \approx 3.07$. 

$$\mathbb{E}(\sqrt{x}) = \int_0^{\infty} \sqrt{x} \frac{1}{\Gamma(\alpha)\beta^{\alpha}}x^{\alpha - 1}e^{-x/\beta}dx$$
$$=\frac{\Gamma(\alpha+0.5)\beta^{0.5}}{\Gamma(\alpha)}\int_0^{\infty} \frac{1}{\Gamma(\alpha+0.5)\beta^{\alpha + 0.5}}x^{\alpha + 0.5 - 1}e^{-x/\beta}dx$$
$$=\frac{\Gamma(\alpha+0.5)\beta^{0.5}}{\Gamma(\alpha)}=\frac{\Gamma(2.8)\sqrt{4.1}}{\Gamma(2.3)}\approx 2.91$$

\subsection*{b)}

The function $g(x) = \ln{x}$ is concave and therefore:

$$g(\mathbb{E}(x)) \geq \mathbb{E}(g(x))$$

Therefore $g(\mathbb{E}(x))$ is an upper bound for $\mathbb{E}(g(x))$. As $\mathbb{E}(x) = 2.3(4.1) = 9.43$ our upper bound is $\ln{9.43} \approx 2.24$. 

$$\mathbb{E}(\ln{x}) = \int_0^{\infty} \ln{x} \frac{1}{\Gamma(\alpha)\beta^{\alpha}}x^{\alpha - 1}e^{-x/\beta}dx$$

According to wolfram alpha this is $\approx 2.01$.

\section*{Question 2}

\subsection*{a)}

First if $Y_i = aX_i$ then:

$$f_{Y_i}(y_i) = \frac{1}{\sqrt{2\pi} \sigma}e^{-y_i^2/(2a^2\sigma^2)}|1/a|$$

That is $Y_i \sim \text{N}(0, a^2\sigma^2)$. \newline

So $Y_1 = -\frac{X_1}{\sqrt{2}\sigma}\sim\text{N}(0, 1/2)$ and $Y_2 = \frac{X_2}{\sqrt{2}\sigma}\sim\text{N}(0, 1/2)$. Let $W=Y_1 + Y_2$. Given these are independent we can look at the product of MGF's:

$$M_W(t) = M_{Y_1}(t)M_{Y_2}(t)=e^{t^2/4}e^{t^2/4}=e^{t^2/2}$$

Which means that $W \sim N(0, 1)$. \newline

Now $U = W^2$ so that:

$$f_U(u) = \frac{1}{\sqrt{2\pi}}e^{-u/2} \frac{1}{\sqrt{u}}$$

Which is exactly the pdf of a chi-squared distribution with degrees of freedom $=1$. 

\subsection*{b)}


The MGF for a $Z \sim \chi^2(k)$ distribution is $(1-2t)^{k/2}$. Therefore for a sum of independent $Z_i \sim \chi^2(k_i)$ distributions we have:

$$M_{\sum_i Z_i}(t) = \prod_i (1-2t)^{k_i/2}=(1-2t)^{\sum_i k_i/2}$$

That is: 

$$W=\sum_i Z_i \sim \chi^2\left(\sum_i k_i\right)$$

In our specific case this means that a sum of $m$ independent random variables that are each $\chi^2(1)$ distributed results in a random variable that is $\chi^2(m)$ distributed. 

\subsection*{c)}

Consider:

$$T/(2\sigma^2)=\sum_{i=1}^{n/2}\left(\frac{X_{2n}-X_{2n-1}}{\sqrt{2}\sigma}\right)^2$$


We already know that each of the terms in that summation are distributed as $\chi^2(1)$ from (a). Furthermore they are all independent as none of the $X_j$ are shared between the terms. Therefore we have the sum of a set of independent random variables distributed as $\chi^2(1)$ which given (b) means:

$$T/(2\sigma^2)\sim \chi^2 (n/2)$$

\subsection*{d)}

The mean of a $\chi^2(m)$ distribution is just $m$. So in our case:

$$\mathbb{E}(T) = 2\sigma^2 n/2=\sigma^2n$$

\subsection*{e)}

$$\mathbb{E}(W) = \mathbb{E}(T/n) = \mathbb{E}(T)/n = \sigma^2n/n = \sigma^2$$

if $n>0$. Therefore $W$ is unbiased for $\sigma^2$. 


\section*{Question 3}

\subsection*{a)}

This is just a $\text{Beta}(\theta + 1, 1)$ distribution. Therefore $\mathbb{E}(X)=\frac{\theta + 1}{\theta + 2}$ and $\text{Var}(X) = \frac{\theta + 1}{(\theta + 2)^2(\theta + 3)}$. From class we know that:

$$\mathbb{E}(\bar{X}) = \mathbb{E}(X) = \frac{\theta + 1}{\theta + 2}$$

$$\text{Var}(\bar{X}) = \text{Var}(X)/n = \frac{\theta + 1}{n(\theta + 2)^2(\theta + 3)}$$

$\bar{X}$ is not unbiased for $\theta$ as $\mathbb{E}(\bar{X}) \neq \theta$. 

\subsection*{b)}

$$\mathbb{E}(\bar{X}) = \mathbb{E}(X) = \theta$$

$$\text{Var}(\bar{X}) = \text{Var}(X)/n = \theta/n$$

Here $\bar{X}$ is unbiased for $\theta$ as $\mathbb{E}(\bar{X}) = \theta$. 

\subsection*{c)}

Let's find $k$ real quick. 

$$\int_0^2 kx(x-2)dx=k\left( x^3/3 - x^2\right)\Big|_0^2=-4k/3\rightarrow k = -3/4$$

$$\mathbb{E}(\bar{X}) = \mathbb{E}(X) = \int_0^2 xkx(x-2)dx = k\left(x^4/4 - 2x^3/3 \right)\Big|_0^2=-4k/3=1$$

$$\mathbb{E}(X^2) = \int_0^2 x^2kx(x-2)dx = k\left(x^5/5 - 2x^4/4 \right)\Big|_0^2=-8k/5=6/5$$

$$\text{Var}(\bar{X}) = \text{Var}(X)/n = \frac{1}{n}\left(6/5 -  (1)^2\right)=\frac{1}{5n}$$

\section*{Question 4}

\subsection*{a)}

From class we know that $\bar{X}_n \sim N(\mu_X, \sigma^2 / n)$. The exact same steps demonstrate that $\bar{Y}_m \sim N(\mu_Y, \sigma^2 / m)$. 

\subsection*{b)}

If $Z = \alpha U$ then we know that:

$$f_Z(z) = f_U(Z/\alpha) |1/\alpha| $$

Therefore: $-\bar{Y}_m \sim N(-\mu_Y, \sigma^2 /m)$. \newline


As the $Y_i$ and $X_i$ are independent $\bar{Y}_m $ and $\bar{X}_n$ are independent as well. Therefore:

$$M_{\bar{X}_n  - \bar{Y}_m}(t) = e^{\mu_X t +\sigma^2 t^2 /n}e^{-\mu_Y t +\sigma^2 t^2 /n}= e^{(\mu_X-\mu_Y) t +2\sigma^2 t^2 /n}$$

Therefore $\bar{X}_n  - \bar{Y}_m \sim N(\mu_X - \mu_Y, 2\sigma^2 / n)$

\subsection*{c)}

From class we know that:

$$\frac{(n-1)S_X^2}{\sigma^2}\sim \chi^2(n-1)$$
$$\frac{(m-1)S_Y^2}{\sigma^2}\sim \chi^2(m-1)$$

and from earlier in this homework we know the sum of two $\chi^2(q)$ independent random variables just results in a distribution that is also $\chi^2$ but with the sum of degrees of freedom. That is:

$$\frac{(n-1)S_X^2}{\sigma^2} + \frac{(m-1)S_Y^2}{\sigma^2} \sim \chi^2(m + n - 2)$$


\section*{Question 5}

Let $Z=T^2$ then:

$$f_Z(z) = f_T(\sqrt{z})\frac{1}{2\sqrt{z}}=k\left(1+\frac{1}{p}z\right)^{-(p+1)/2}\frac{1}{2\sqrt{z}}=\frac{k}{2}\frac{z^{1/2-1}}{(1+(1/p)z)^{(1+p)/2}}$$

Which is the same functional form as $F(1,p)$. Therefore knowing that $k/2$ will result in the normalization of this function into a pdf $T^2\sim F(1,p)$. 

\section*{Question 6}

By definition we know that if $X\sim F(p, q)$ then there exists a $U\sim \chi^2 (p)$ and $V\sim \chi^2 (q)$ s.t.

$$X=\frac{U/p}{V/q}$$

so that $X\sim F(p,q)$. \newline

In that case: $Y=1/X=\frac{V/q}{U/p}$. Therefore it follows that $Y\sim F(q,p)$. 

\end{document}